% Capitulo: De usuários a tokens: custo e consumo
\chapter{De usuários a tokens: custo e consumo}
\label{cap:custo-consumo}

Durante muito tempo, o custo de uma ferramenta de dados era uma linha fixa no orçamento: um
valor por usuário, negociado na renovação do contrato. O uso, alto ou baixo, não mudava a
fatura. Essa lógica está mudando. Os armazéns de dados na nuvem passaram a cobrar pelo
processamento, as ferramentas de BI passaram a oferecer modelos por capacidade ou por
crédito, e os modelos de linguagem cobram por token, a unidade de texto que entra e sai de
cada pergunta. Com isso, gerir o consumo virou uma preocupação de todo gerente de dados. Este
capítulo trata dessa mudança, do que ela exige de quem lidera um time e de como os tokens
entram na conta.

\section{O modelo por usuário}
\label{sec:modelo-usuario}

No modelo tradicional, paga-se uma licença por pessoa com acesso à ferramenta, muitas vezes
em níveis diferentes: quem cria conteúdo paga mais, quem só visualiza paga menos. O custo é
previsível, e é negociado de tempos em tempos, quase sempre por compras e TI.

A alavanca de controle do gerente, nesse modelo, é o acesso: quem precisa de licença e de
qual nível. O desperdício típico também é de acesso: licenças pagas para pessoas que saíram da
área, que nunca usaram a ferramenta ou que precisariam só do nível mais barato. É um
desperdício real, mas limitado e fácil de enxergar numa planilha de usuários.

O que o modelo por usuário não faz é ligar o custo ao uso. Um painel que atualiza de hora em
hora custa o mesmo que um que atualiza uma vez por dia; uma consulta que lê a base inteira custa
o mesmo que uma bem escrita. Para o orçamento, a eficiência do trabalho técnico é invisível.

\section{Créditos, consumo e tokens}
\label{sec:modelo-consumo}

No modelo por consumo, paga-se pelo que foi processado: horas de computação, volume de dados
lido por consulta, atualizações executadas, tokens de entrada e de saída de um modelo de
linguagem. A Tabela~\ref{tab:usuario-consumo} compara os dois modelos.

\begin{table}[htbp]
  \centering
  \caption{Cobrança por usuário e por consumo.}
  \label{tab:usuario-consumo}
  \small
  \begin{tabularx}{\textwidth}{@{}>{\raggedright\arraybackslash}p{0.2\textwidth}
      >{\raggedright\arraybackslash}X>{\raggedright\arraybackslash}X@{}}
    \toprule
    & Por usuário & Por consumo (créditos, tokens) \\
    \midrule
    O que se paga & Uma licença por pessoa com acesso & O que foi processado: consultas, atualizações, tokens \\
    \addlinespace
    Previsibilidade & Alta, custo fixo por período & Baixa, varia com o uso \\
    \addlinespace
    Alavanca de controle & Quem tem acesso e com qual nível & O que roda, com que frequência e com qual modelo \\
    \addlinespace
    Desperdício típico & Licença de quem não usa & Atualização ou consulta que ninguém aproveita \\
    \addlinespace
    Quem acompanha & Compras e TI, na renovação & O gerente do time, de forma contínua \\
    \bottomrule
  \end{tabularx}
  \fonte{elaborado pelo autor.}
\end{table}

A mudança começou com os armazéns de dados na nuvem, que separaram armazenamento de
processamento e passaram a cobrar o processamento pelo uso (\veja{sec:dw-nuvem}). A IA
acelerou a tendência, porque cada pergunta feita a um modelo tem um custo em tokens, e cada
agente que executa uma tarefa faz dezenas ou centenas de chamadas ao modelo. E ferramentas de
BI e de dados que antes cobravam só por usuário passaram a oferecer, ou a impor, modelos por
capacidade ou por crédito. \pendente{conferir e citar exemplos atuais de ferramentas de BI
com cobrança por crédito ou capacidade}

O modelo por consumo tem vantagens reais. Paga-se pelo uso efetivo, um time pequeno começa
barato e cresce conforme precisa, e a eficiência técnica passa a aparecer na fatura, o que dá
ao time um incentivo concreto para fazer bem feito. O risco é a imprevisibilidade: um processo
mal configurado, um agente em laço ou uma consulta pesada agendada por engano podem consumir o
orçamento do mês em poucos dias, e só aparecer na fatura depois.

\antesdepois{O custo da ferramenta era uma linha fixa no orçamento. O gerente controlava
  quem tinha licença, e o uso, alto ou baixo, não mudava a fatura.}
{O custo acompanha o uso. O gerente controla o que roda, com que frequência e com qual
  modelo, e acompanha o consumo como acompanha qualquer outro indicador do time.}

\section{O gerente como gestor de consumo}
\label{sec:gestao-consumo}

No modelo por consumo, o gerente de dados passa a ser, também, gestor de um orçamento
variável. Isso não é exclusivo de dados: a disciplina de gestão de custos de nuvem, conhecida
como FinOps, propõe exatamente isso para qualquer área que consome recursos de nuvem, com três
ideias centrais: dar visibilidade ao custo, tornar cada time responsável pelo próprio consumo
e otimizar de forma contínua, e não em campanhas \parencite{storment2023,finops2026}.

Para um time de dados, algumas práticas traduzem essas ideias para o dia a dia:

\begin{itemize}
  \item \textbf{Orçamento por time e por produto}, com alertas quando o consumo se aproxima do
    limite, e não só depois que a fatura chega.
  \item \textbf{Custo por produto de dados}, a partir das métricas de metadados do capítulo
    anterior (\veja{sec:metricas-uso}): quanto custa manter cada painel, cada pipeline, cada
    assistente.
  \item \textbf{Regras técnicas que evitam desperdício}, verificadas automaticamente na
    integração contínua: não ler todas as colunas de uma tabela em modelos de consumo, não
    fazer junções sem condição, preferir cargas incrementais (\veja{sec:ci}).
  \item \textbf{Ambientes dimensionados para cada uso}, com o processamento de
    desenvolvimento separado e menor que o de produção, e desligado quando ninguém usa.
  \item \textbf{Custo visível para quem pede.} Quando o stakeholder vê quanto custa atualizar
    o painel de hora em hora, a conversa sobre frequência muda de natureza
    (\veja{sec:papeis-frequencia}).
\end{itemize}

A última prática é a que mais muda a relação com o negócio. No modelo por usuário, pedir mais
uma atualização ou mais um painel parecia gratuito. No modelo por consumo, cada pedido tem
um custo que pode ser mostrado, e a decisão passa a ser de valor: vale pagar isso por esse
benefício?

\rascunho{
  Relato: como eu acompanho o consumo hoje (que indicadores olho, com que frequência, que
  alertas existem) e um caso de otimização que reduziu o custo de forma relevante. Sem nomes
  e com valores só em ordem de grandeza.
}

\section{Tokens no trabalho com dados}
\label{sec:tokens-dados}

Os tokens aparecem em vários pontos do trabalho com dados: no assistente que ajuda a escrever
consultas, na pergunta em linguagem natural feita sobre um painel, no agente que executa uma
tarefa inteira, na documentação gerada automaticamente, no motor de conciliação que pede uma
sugestão à IA (\veja{sec:motor-conciliacao}). Em cada caso, o custo depende do tamanho do que
entra (o pedido e o contexto) e do que sai (a resposta), e do modelo usado.

Três fatores costumam fazer o custo crescer mais do que o necessário. O primeiro é o contexto
grande: mandar a tabela inteira, a documentação inteira ou o histórico inteiro a cada
pergunta. Além de caro, nem sempre melhora a resposta (\veja{sec:modelo-so-sabe}). O segundo
são as perguntas repetidas sem reaproveitamento: a mesma pergunta, feita por dez pessoas,
processada dez vezes. O terceiro é usar o modelo mais caro para tudo, inclusive para tarefas
simples que um modelo menor faria igualmente bem.

As práticas que ajudam seguem a mesma lógica do resto do capítulo. Pré-calcular o que é
perguntado sempre, em vez de responder do zero a cada vez. Escolher o modelo pelo tamanho da
tarefa. Limitar o contexto ao necessário, com recuperação de trechos relevantes em vez de
documentos inteiros. E medir tokens por resposta útil, do mesmo jeito que se mede créditos por
acesso num painel.

Vale, por fim, a mesma pergunta do capítulo anterior, agora aplicada à IA: não adianta
responder em segundos, com um modelo caro, uma pergunta que poderia ser atendida por um
relatório semanal. A frequência e a forma de entrega certas continuam sendo decisões de
negócio, e são elas que separam o uso da IA que gera valor do que só gera fatura.

\begin{leituras}
  \leitura{storment2023}{A referência sobre FinOps: como dar visibilidade ao custo de
    nuvem, distribuir a responsabilidade entre os times e otimizar sem travar o trabalho.}
  \leitura{finops2026}{O modelo mantido pela comunidade de FinOps, atualizado com frequência
    e com material específico sobre custos de IA.}
  \leitura{dageville2016}{Para entender, na fonte, por que o custo dos armazéns de dados na
    nuvem passou a acompanhar o processamento.}
\end{leituras}
